\documentclass[10pt]{article}

\usepackage[utf8]{inputenc}

\usepackage[T1]{fontenc}

\usepackage{amsmath}

\usepackage{amsfonts}

\usepackage{amssymb}

\usepackage[version=4]{mhchem}

\usepackage{stmaryrd}

\usepackage{graphicx}

\usepackage[export]{adjustbox}

\graphicspath{ {./images/} }

\usepackage{bbold}



\def\AA{\mathring{\mathrm{A}}}



\begin{document}

Jum.: [1] F r o b e n i u s F., «J. reine und angew. Math.», 1877, Bd 82, S. 230-315; [2] X'а р т м а н Ф., Обыкновенные дифференциальные уравнения, пер. с англ., М , 1970



ФРОБЕНИУСА ФOPMУJА - формула, выра шая отношение обобшенного определителя Вандермонда к обычному (см. Вандермонда определитель) через степенные суммы. В качестве коэффициентов в Ф. ф. участвуют характеры представлений симметрической apynnbl.



Пусть $x_{1}, \ldots, x_{n}$ - независимые переменные. Для любого набора $\lambda=\left(\lambda_{1}, \ldots, \lambda_{n}\right)$ неотрицательных целых чисел, удовлетворяющего условию $\lambda_{1} \geqslant \lambda_{2} \geqslant \ldots \geqslant \lambda_{n}$, IIусть



\begin{center}

\includegraphics[max width=\textwidth]{2023_07_09_f11aefc4a610aea3782ag-1}

\end{center}



так что $W_{0}$ есть обычный огределитель Вандермонда. И пусть $\sum \lambda_{i}=m$; тогда набор $\lambda$ после выкидывания нулей можно рассматривать как разбиение числа $m$. Рассматривается соответствующее неприводимое представление $T_{\lambda}$ грушпы $S_{m}$. Для любого разбиения $\mu=\left(\mu_{1}, \ldots, \mu_{r}\right)$ тісла $m$ через $a_{\lambda \mu}$ обозначается знатение характера представления $T_{\lambda}$ на классе сопряженных әлементов группы $S_{m}$, определяемом разбиением $\mu$, и через $c_{\mu}$-порядок централизатора любой подстановки из этого класса. Пусть $s_{\mu}=s_{\mu_{1}} s_{\mu_{2}} \ldots s_{\mu_{r}}$, где $s_{k}=x_{1}^{k}+\ldots+x_{n}^{k}$. Тогда



\[

\frac{W_{\lambda}}{W_{0}}=\sum_{\mu} a_{\lambda \mu} c_{\mu}^{-1} s_{\mu}

\]



где сумма берется по всем (неупорядоченным) раабисниям числа $m$. При этом, если разбиение $\mu$ содержит $k_{1}$ единиц, $k_{2}$ двоек и т. Д., то



\[

c_{\mu}=k_{1} ! k_{2} ! \ldots 1^{k_{1}} 2^{k_{2}} \ldots

\]



Если $n \geqslant m$, то Ф. ф. может быть преобразована к виду



\[

\sum_{\lambda} a_{\lambda \mu} W_{\lambda}=s_{\mu} W_{0}

\]



где сумма берется по всем разбиениям числа $m$ (дополненным надлежащим числом нулей). Последняя формула может быть использована для вычисления характеров симметріч. группы. А именно, $a_{\lambda \mu}$ есть коэффициент при $x_{1}^{\lambda_{1}+n-1} x_{2}^{\lambda_{2}+n-2} \ldots x_{n}^{\lambda_{n}}$ в многочлене $s_{\mu} W_{0} \cdot$ $\begin{aligned} & \text { Лum.: [1] М у р н а г а н Ф. Д., Теория представлений } \\ & \text { групп, пер. с англ, М., 1950. } \text { Э. Б. Винбера. }\end{aligned}$ ФРОБЕНИУСА ' ЭНДОМОРФИЗМ - эндоморфизм $\varphi: X \longrightarrow X$ схемы $X$ над конечным полем $\mathrm{F}_{q}$ из $q$ элементов такой, что $\varphi$-тождественное отображение топологич. пространства $X$, а отображение структурного пучка $\varphi^{*}: G_{X} \rightarrow \mathcal{G}_{X}$ совпадает с возведением в степень $q$. Ф. э. является чисто несепарабельным морфизмом и имеет нулевой дифференциал. Для аффинного многообразия $X \subset A^{n}$, определенного над $\mathrm{F}_{q}$, Ф. э. $\varphi$ переводит точку $\left(x_{1}, \ldots, x_{n}\right)$ в точку $\left(x_{1}^{q}, \ldots, x_{n}^{q}\right)$.



Число геометрич. точек схемы $X$, определенных над $\mathrm{F}_{q}$, совпадает с числом неподвижных точек Ф. э. $\varphi$, что позволяет использовать для определения числа таких точек Лефшеца бормулу.



Лum.: [1] X а p т с х о p H P., Алгебраическая геометрия, пер. с англ., М., 1981. Л. В. Кузъмин.



ФРОБЕНИУСОВА АЛГЕБРА - конечномерная алгебра $R$ над полем $P$ такая, что левые $R$-модули $R$ и $\operatorname{Hom}_{P}(R, P)$ изоморфны. На языке представлений это означает әквивалентность правого и левого регулярных представлений. Всякая групшовая алгебра конеч- ной группы над полем является Ф. а. Каждая Ф. а. является квазибробениусовым кольцом. Обратное утверждение неверно. Эквивалентны следующие свойства конечномерной $P$-алгебры $R$ :



\begin{enumerate}

  \item $R-\Phi$. а.;



  \item существует такая невырожденная билинейная форма $f: R \times R \rightarrow P$, что $f(a b, c)=f(a, b c)$ для любых $a, b$, $c \in R$



  \item если $L$ - левый, а $H$ - правый пдеалы алгебры $R$, то (см. Аннулятор)



\end{enumerate}



\[

B_{l}\left(3_{r}(L)\right)=L, \quad 3_{r}\left(3_{l}(H)\right)=H \text {, }

\]



$\operatorname{dim}_{P} \AA_{r}(L)+\operatorname{dim}_{P} L=\operatorname{dim}_{P} R=\operatorname{dim}_{P} \bigcap_{l}(H)+\operatorname{dim}_{P} H$.



Ф. а., по существу, появились впервые в работах Г. Фробениуса [3].



Лuт.: [1] К э р т и с Ч., Р а й н е р И., Теория представлений нонечных групп и ассоциативных алгебр, пер. с англ., M., 1969; [2] Ф е и с К., Алгебра: кольца, модули и категории, "Sitzungsber. Königl. Preuss. Akad. Wiss.», 1903 , N 24, S. 504-, 507, 634-45.



ФРОММЕРА МЕТОД - метод исследования особых точек автономной системы обыкновенных диффференциальных уравнений 2-го порядка



\[

\dot{p}=f(p), \quad p=(x, y), f=(X, Y): G \longrightarrow \mathbb{R}^{2},

\]



где функция $f$ аналитическая или достаточно гладкая в области $G$.



Пусть $O=(0,0)$ - особая точка системы (1), т. е. $f(O)=0$, а $X, Y$ - аналитические в точке $O$ функции, не имеющие общего аналитического исчезающего в $O$ множителя. Ф. м. позволяет выявить все $T O-к \mathrm{p}$ ив ы е системы (1) - полутраектории, әтой системы, примыкающие к $O$ по определенным направлениям.



\includegraphics[max width=\textwidth, center]{2023_07_09_f11aefc4a610aea3782ag-1(1)}

$x=0$, является $O-\kappa$ р и в о й уравнения



\[

y^{\prime}=\frac{Y(x, y)}{X(x, y)}

\]



(т. е. представима вблизи точки $O$ в виде



\[

y=\dot{\varphi}(x), \quad \varphi(x) \longrightarrow 0 \quad \text { при } \quad x \longrightarrow 0

\]



где $\varphi: I \rightarrow \mathbb{R}$ - решение уравнения $(2), I=(0, \delta)$ или $(-\delta, 0), \delta>0, \varphi(x) \equiv 0$ или $\varphi(x) \neq 0$ для любого $x \in I)$ и наоборот.



Пусть уравнение (2) рассматривается сначала в области $x>0$. Если оно представляет собой п р о с т о е у р а в е ение Бенд ик с он а, т. е. удовлетворяет условиям



\[

X(x, y) \equiv x^{h}, \quad h \geqslant 1, \quad Y_{y}^{\prime}(0,0)=a \neq 0,

\]



то для него в области $x>0$ существует единственная $O$-кривая при $a<0 ;$ область $x>0, x^{2}+y^{2}<r^{2}$, где $r-$ достаточно малое положительное число, представляет собой параболич. сектор при $a>0$. В противном случае для выявления $O$-кривых уравнения (2) в области $x>0$ применяют Ф. м. Основой для его применения является тот факт, что каждая $O$-кривая (3) уравнения $(2), \varphi(x) \neq 0$, обладает в точке $O$ вполне определенной асимптотикой, а именно, будучи представлена в виде



\[

y=x^{v(x)} \text { sign } \varphi(x),

\]



она допускает конечный или бесконечный предел



\[

v=\lim _{x \rightarrow 0} v(x)=\lim _{x \rightarrow 0} \frac{\ln |\varphi(x)|}{\ln x} \in[0,+\infty],

\]



к-рый наз. ее $п о р$ я д к о м к р и в и 3 ны в точке $O$, a при $v \in(0,+\infty)$ допускает еще и конечный или бесконечный предел



\[

\gamma=\lim _{x \rightarrow 0} \varphi(x) x^{-v} \in[-\infty,+\infty]

\]



к-рый наз, ее ме р о й к р и ви з ны в точке $O$. При этом $O$-кривой $y=0, \quad x \in(0, \delta)$, приписывается порядок кривизны $v=+\infty$.





\end{document}